% Gerado por ferramentas/gerar.py; não edite à mão.
\section*{Exercícios}\markright{Exercícios}
\addcontentsline{toc}{section}{Exercícios}
Os pontos antes de cada número indicam a dificuldade: \nivel{1} fácil, \nivel{2} médio, \nivel{3} difícil e \nivel{4} desafio. Os exercícios marcados com \estrela{} têm solução completa, com animação do gráfico, no repositório do livro: \url{https://github.com/AleffAzevedo/calculo-2-pt/blob/main/exercicios/cap2/README.md}. As respostas estão no fim do capítulo.

\begin{itemize}[leftmargin=5.2em,labelsep=0.6em,itemsep=4pt]
\item[\exnum{1}{2.1}] Calcule $\displaystyle\lim_{n\to\infty}\frac{2n^2+3n}{5n^2-1}$.
\item[\exnum{1}{2.2}] Calcule $\displaystyle\lim_{n\to\infty}n\sen\frac1n$.
\item[\exnum{1}{2.3}] Calcule $\displaystyle\lim_{n\to\infty}\Big(1+\frac2n\Big)^{3n}$.
\item[\exnum{1}{2.4}] Calcule $\displaystyle\lim_{n\to\infty}\big(\sqrt{n^2+n}-n\big)$.
\item[\exnum{1}{2.5}] A sequência $a_n=(-1)^n\dfrac n{n+1}$ converge?
\item[\exnum{1}{2.6}] Calcule $\displaystyle\lim_{n\to\infty}\frac{n^2}{2^n}$.
\item[\exnum{1}{2.7}] Calcule $\displaystyle\sum_{n=0}^\infty\Big(-\frac23\Big)^n$.
\item[\exnum{1}{2.8}] Calcule $\displaystyle\sum_{n=1}^\infty\frac{5\cdot3^n}{4^{n+1}}$.
\item[\exnum{1}{2.9}] Escreva $1{,}2\overline{3}=1{,}2333\ldots$ como fração usando uma série geométrica.
\item[\exnum{1}{2.10}] Calcule $\displaystyle\sum_{n=1}^\infty\frac1{(n+1)(n+2)}$ (telescópica).
\item[\exnum{1}{2.11}] Mostre que $\displaystyle\sum_{n=1}^\infty\ln\frac n{n+1}$ diverge calculando as somas parciais.
\item[\exnum{1}{2.12}] Use o teste da divergência em $\sum\frac{n^2}{n^2+1}$ e em $\sum\cos\frac1n$.
\item[\exnum{1}{2.13}] Classifique: $\sum\frac1{n^{1{,}01}}$, $\sum\frac1{\sqrt n}$, $\sum n^{-\pi}$.
\item[\exnum{1}{2.14}] Use o teste da integral para decidir sobre $\displaystyle\sum_{n=2}^\infty\frac1{n\ln n}$.
\item[\exnum{1}{2.15}] Use o teste da integral em $\sum n\,e^{-n^2}$.
\item[\exnum{1}{2.16}] Use comparação direta: $\sum\frac1{n^3+n}$ e $\sum\frac{2+\sen n}{n}$.
\item[\exnum{1}{2.17}] Use comparação no limite em $\displaystyle\sum\frac{n+1}{n^3+2}$.
\item[\exnum{1}{2.18}] Classifique $\displaystyle\sum\frac{(-1)^n}{\sqrt n}$ (absoluta, condicional ou divergente).
\item[\exnum{1}{2.19}] Use o teste da razão em $\displaystyle\sum\frac n{3^n}$.
\item[\exnum{1}{2.20}] Use o teste da razão em $\displaystyle\sum\frac{3^n}{n^2}$.
\item[\exnum{2}{2.21}] \estrela\ \textbf{Raízes encaixadas.} Seja $a_1=\sqrt2$ e $a_{n+1}=\sqrt{2+a_n}$, isto é, $\sqrt2,\ \sqrt{2+\sqrt2},\ \sqrt{2+\sqrt{2+\sqrt2}},\dots$
Mostre que a sequência converge e calcule o limite.
\item[\exnum{2}{2.22}] \estrela\ \textbf{A bola que quica.} Uma bola cai de $10$ m e, a cada quique, sobe a $\frac34$ da altura anterior. Qual a distância
total percorrida? E o tempo total, se a queda de uma altura $h$ leva $\sqrt{2h/g}$ ($g=9{,}8$)?
\item[\exnum{2}{2.23}] \estrela\ \textbf{A lentidão da série harmônica.} Use o teste da integral para mostrar que $\ln(n+1)\le H_n\le1+\ln n$, em que $H_n=1+\frac12+\dots+\frac1n$.
Quantos termos são necessários para $H_n$ passar de $10$?
\item[\exnum{2}{2.24}] \estrela\ \textbf{Telescópica com salto.} Calcule $\displaystyle\sum_{n=1}^\infty\frac1{n(n+2)}$.
\item[\exnum{2}{2.25}] \estrela\ \textbf{Estimando $\sum 1/n^2$.} Calcule $S_{10}$ para $\sum\frac1{n^2}$ e use as cotas do resto pelo teste da integral para obter um intervalo para a soma.
Compare com o valor exato $\frac{\pi^2}6$.
\item[\exnum{2}{2.26}] \estrela\ \textbf{Uma alternada que converge depressa.} A série $\displaystyle\sum_{n=1}^\infty\frac{(-1)^{n+1}}{n!}$ vale $1-\frac1e$. Quantos termos garantem erro menor que $10^{-4}$?
Compare a soma parcial com o valor exato.
\item[\exnum{2}{2.27}] Calcule $\displaystyle\sum_{n=1}^\infty\frac{2^n+3^n}{6^n}$.
\item[\exnum{2}{2.28}] Para quais $x$ a série $\sum_{n=0}^\infty(x-1)^n$ converge? Qual a soma?
\item[\exnum{2}{2.29}] Calcule $\displaystyle\lim_{n\to\infty}\sqrt[n]n$.
\item[\exnum{2}{2.30}] Mostre que $\dfrac{n!}{n^n}\le\dfrac1n$ e conclua o limite.
\item[\exnum{2}{2.31}] Calcule $\displaystyle\sum_{n=2}^\infty\frac1{n^2-1}$.
\item[\exnum{2}{2.32}] Decida: $\displaystyle\sum\frac{\arctg n}{n^2}$.
\item[\exnum{2}{2.33}] Decida: $\displaystyle\sum_{n=2}^\infty\frac1{n(\ln n)^2}$.
\item[\exnum{2}{2.34}] Decida: $\displaystyle\sum\frac{n^2+1}{n^4-n+3}$.
\item[\exnum{2}{2.35}] Decida: $\displaystyle\sum\sen\frac1n$.
\item[\exnum{2}{2.36}] Classifique $\displaystyle\sum\frac{(-1)^nn}{n^2+1}$.
\item[\exnum{2}{2.37}] Decida: $\displaystyle\sum\frac{(-1)^nn!}{10^n}$.
\item[\exnum{2}{2.38}] Decida: $\displaystyle\sum\frac{n^n}{n!}$.
\item[\exnum{2}{2.39}] Decida com o teste da raiz: $\displaystyle\sum\Big(\frac n{n+1}\Big)^{n^2}$.
\item[\exnum{2}{2.40}] Decida: $\displaystyle\sum\frac1{3^n-n}$.
\item[\exnum{3}{2.41}] \estrela\ \textbf{O floco de neve de Koch.} Partindo de um triângulo equilátero de área $A_0$, a cada etapa cada lado é dividido em três e, sobre o terço
do meio, constrói-se um novo triângulo equilátero para fora. Mostre que o perímetro tende a infinito, mas a área tende a $\frac85A_0$.
\item[\exnum{3}{2.42}] \estrela\ \textbf{Reordenar muda a soma.} Sabendo que $1-\frac12+\frac13-\frac14+\dots=\ln2$, mostre que a reordenação
$1+\frac13-\frac12+\frac15+\frac17-\frac14+\cdots$ (dois positivos, um negativo) tem soma $\frac32\ln2$.
\item[\exnum{3}{2.43}] \estrela\ \textbf{Empilhando blocos.} Blocos de comprimento $1$ são empilhados na borda de uma mesa. Colocando o centro de massa dos $k$ blocos de cima
sobre a borda do bloco $k+1$, o $k$-ésimo bloco (de cima para baixo) avança $\frac1{2k}$. Qual o balanço total com $n$ blocos?
Quantos blocos são precisos para o bloco do topo ficar inteiramente fora da mesa? E para um balanço de $3$?
\item[\exnum{3}{2.44}] \estrela\ \textbf{O conjunto de Cantor.} Do intervalo $[0,1]$ retira-se o terço do meio $(\frac13,\frac23)$; de cada um dos dois intervalos restantes,
retira-se o terço do meio; e assim por diante. Qual o comprimento total retirado?
\item[\exnum{3}{2.45}] \estrela\ \textbf{Somando $\sum n/2^n$.} Mostre que $\displaystyle\sum_{n=1}^\infty\frac n{2^n}$ converge e calcule a soma usando $S-\frac12S$.
Faça o mesmo para $\displaystyle\sum_{n=1}^\infty\frac{n^2}{2^n}$.
\item[\exnum{3}{2.46}] Decida: $\displaystyle\sum_{n=3}^\infty\frac1{n\ln n\,\ln(\ln n)}$.
\item[\exnum{3}{2.47}] Para quais $p$ a série $\displaystyle\sum_{n=2}^\infty\frac1{n(\ln n)^p}$ converge?
\item[\exnum{3}{2.48}] Decida: $\displaystyle\sum\frac{\ln n}{n^{1{,}5}}$.
\item[\exnum{3}{2.49}] Decida: $\displaystyle\sum\frac1{n^{1+1/n}}$.
\item[\exnum{3}{2.50}] Decida: $\sum\big(\sqrt{n+1}-\sqrt n\big)$ e $\sum\dfrac{\sqrt{n+1}-\sqrt n}{n}$.
\item[\exnum{3}{2.51}] Classifique $\displaystyle\sum_{n=2}^\infty\frac{(-1)^n}{\ln n}$.
\item[\exnum{3}{2.52}] Classifique $\displaystyle\sum(-1)^n\Big(1-\cos\frac1n\Big)$.
\item[\exnum{3}{2.53}] Decida: $\displaystyle\sum\frac{n^{100}}{1{,}01^n}$.
\item[\exnum{3}{2.54}] Decida: $\displaystyle\sum\frac{(n!)^2}{(2n)!}$.
\item[\exnum{3}{2.55}] Decida: $\displaystyle\sum\frac{1\cdot3\cdot5\cdots(2n-1)}{2\cdot4\cdot6\cdots(2n)}$. Por que o teste da razão não basta?
\item[\exnum{3}{2.56}] Seja $a_1=1$ e $a_{n+1}=\dfrac n{2n+1}a_n$. A série $\sum a_n$ converge?
\item[\exnum{3}{2.57}] Calcule $\displaystyle\sum_{n=1}^\infty\frac1{n(n+1)(n+2)}$.
\item[\exnum{3}{2.58}] Mostre que $\sum\frac{(-1)^n}{n^p}$ converge para todo $p>0$ e que a convergência é absoluta se, e só se, $p>1$.
\item[\exnum{3}{2.59}] Quantos termos de $\sum\frac{(-1)^{n+1}}{n^3}$ garantem erro menor que $0{,}0005$? Calcule a aproximação.
\item[\exnum{3}{2.60}] Seja $a_1=\tfrac12$ e $a_{n+1}=a_n(2-a_n)$. Mostre que a sequência é crescente e limitada por $1$ e calcule o limite.
\item[\exnum{4}{2.61}] Mostre que $\displaystyle\sum_{n=1}^k\frac1{n^2}<2-\frac1k$ usando $\frac1{n^2}<\frac1{n(n-1)}$ e conclua que $\sum\frac1{n^2}\le2$.
\item[\exnum{4}{2.62}] Seja $c_n=1+\frac12+\dots+\frac1n-\ln n$. Mostre que $\{c_n\}$ é decrescente e limitada inferiormente, logo converge (o limite é a constante de Euler $\gamma\approx0{,}5772$).
\item[\exnum{4}{2.63}] Usando $e=\sum_{n=0}^\infty\frac1{n!}$, prove que $e$ é irracional. (Suponha $e=\frac pq$ e multiplique por $q!$.)
\item[\exnum{4}{2.64}] (Condensação de Cauchy) Se $a_n\ge0$ é decrescente, mostre que $\sum a_n$ converge se, e só se, $\sum2^ka_{2^k}$ converge. Use isso para refazer as séries $p$.
\item[\exnum{4}{2.65}] Se $a_n>0$ e $\sum a_n$ converge, prove que $\sum a_n^2$ converge. Mostre com um exemplo que $\sum\sqrt{a_n}$ pode divergir.
\end{itemize}
\subsection*{Respostas}\markright{Respostas}
{\small\begin{multicols}{2}\raggedright
\begin{itemize}[leftmargin=2.7em,labelsep=0.5em,itemsep=1pt]
\item[\textbf{2.1}] $\tfrac25$.
\item[\textbf{2.2}] $1$ (com $t=\frac1n$, $\frac{\sen t}t\to1$).
\item[\textbf{2.3}] $e^6$.
\item[\textbf{2.4}] $\tfrac12$ (multiplique pelo conjugado).
\item[\textbf{2.5}] Não: os termos pares tendem a $1$ e os ímpares a $-1$.
\item[\textbf{2.6}] $0$ (L'Hôpital duas vezes em $\frac{x^2}{2^x}$).
\item[\textbf{2.7}] $\tfrac35$.
\item[\textbf{2.8}] $\tfrac{15}4$.
\item[\textbf{2.9}] $1{,}2+\tfrac{0{,}03}{1-0{,}1}=\tfrac{37}{30}$.
\item[\textbf{2.10}] $\tfrac12$.
\item[\textbf{2.11}] $S_k=-\ln(k+1)\to-\infty$.
\item[\textbf{2.12}] Ambas divergem (termos $\to1$).
\item[\textbf{2.13}] Converge; diverge; converge.
\item[\textbf{2.14}] $\int_2^\infty\tfrac{dx}{x\ln x}=\infty$: diverge.
\item[\textbf{2.15}] $\int_1^\infty xe^{-x^2}dx=\tfrac1{2e}$: converge.
\item[\textbf{2.16}] Converge ($<\tfrac1{n^3}$); diverge ($\ge\tfrac1n$).
\item[\textbf{2.17}] Com $\tfrac1{n^2}$, $L=1$: converge.
\item[\textbf{2.18}] Condicionalmente convergente.
\item[\textbf{2.19}] $\rho=\tfrac13$: converge.
\item[\textbf{2.20}] $\rho=3$: diverge.
\item[\textbf{2.21}] Crescente e limitada por $2$; o limite é $2$.
\item[\textbf{2.22}] $70$ m; cerca de $19{,}9$ s.
\item[\textbf{2.23}] $12\,367$ termos.
\item[\textbf{2.24}] $\tfrac34$.
\item[\textbf{2.25}] $S_{10}\approx1{,}5498$; $1{,}6407\le S\le1{,}6498$; $\frac{\pi^2}6\approx1{,}6449$.
\item[\textbf{2.26}] $N=7$; $S_7\approx0{,}632143$ contra $0{,}632121$.
\item[\textbf{2.27}] $\tfrac12+1=\tfrac32$.
\item[\textbf{2.28}] $0<x<2$; soma $\tfrac1{2-x}$.
\item[\textbf{2.29}] $1$ ($\ln\sqrt[n]n=\tfrac{\ln n}n\to0$).
\item[\textbf{2.30}] $\tfrac{n!}{n^n}=\tfrac1n\cdot\tfrac2n\cdots\tfrac nn\le\tfrac1n\to0$.
\item[\textbf{2.31}] $\tfrac12\big(\tfrac1{n-1}-\tfrac1{n+1}\big)$: soma $\tfrac34$.
\item[\textbf{2.32}] Converge ($\le\tfrac{\pi/2}{n^2}$).
\item[\textbf{2.33}] Converge ($\int_2^\infty=\tfrac1{\ln2}$).
\item[\textbf{2.34}] Converge (como $\tfrac1{n^2}$).
\item[\textbf{2.35}] Diverge (comparação no limite com $\tfrac1n$).
\item[\textbf{2.36}] Condicionalmente convergente.
\item[\textbf{2.37}] Diverge: os termos crescem sem limite.
\item[\textbf{2.38}] $\rho=e>1$: diverge.
\item[\textbf{2.39}] $\sqrt[n]{a_n}=\big(\tfrac n{n+1}\big)^n\to\tfrac1e<1$: converge.
\item[\textbf{2.40}] Converge (comparação no limite com $\tfrac1{3^n}$).
\item[\textbf{2.41}] Perímetro $3\ell(4/3)^n\to\infty$; área $A_0\big(1+\sum_{n\ge1}\frac{3\cdot4^{n-1}}{9^n}\big)=\frac85A_0$.
\item[\textbf{2.42}] $\tfrac32\ln2\approx1{,}0397$.
\item[\textbf{2.43}] $\frac12H_n$; $4$ blocos; $227$ blocos.
\item[\textbf{2.44}] $1$: o conjunto que sobra tem comprimento total zero.
\item[\textbf{2.45}] $2$ e $6$.
\item[\textbf{2.46}] Diverge (integral: $\ln\ln\ln x\to\infty$).
\item[\textbf{2.47}] $p>1$.
\item[\textbf{2.48}] Converge: $\ln n<n^{0{,}25}$ para $n$ grande, compare com $n^{-1{,}25}$.
\item[\textbf{2.49}] Diverge: $\frac{1/n^{1+1/n}}{1/n}=n^{-1/n}\to1$.
\item[\textbf{2.50}] A primeira diverge ($S_k=\sqrt{k+1}-1$); a segunda converge (termos $\sim\tfrac1{2n^{3/2}}$).
\item[\textbf{2.51}] Condicionalmente convergente ($\tfrac1{\ln n}>\tfrac1n$).
\item[\textbf{2.52}] Absolutamente convergente ($1-\cos\tfrac1n\sim\tfrac1{2n^2}$).
\item[\textbf{2.53}] $\rho=\tfrac1{1{,}01}<1$: converge (a exponencial vence).
\item[\textbf{2.54}] $\rho=\tfrac14$: converge.
\item[\textbf{2.55}] $\rho=1$; mas $a_n\ge\tfrac12\prod_{k=2}^n\tfrac{2k-2}{2k}=\tfrac1{2n}$: diverge por comparação com a harmônica.
\item[\textbf{2.56}] $\rho=\tfrac12$: converge.
\item[\textbf{2.57}] $\tfrac12\big(\tfrac1{n(n+1)}-\tfrac1{(n+1)(n+2)}\big)$: soma $\tfrac14$.
\item[\textbf{2.58}] Leibniz para $p>0$; $\sum\tfrac1{n^p}$ converge sse $p>1$.
\item[\textbf{2.59}] $\tfrac1{(N+1)^3}<0{,}0005$: $N=12$; $S_{12}\approx0{,}9013$ (valor $0{,}9015$).
\item[\textbf{2.60}] $1-a_{n+1}=(1-a_n)^2$; limite $1$.
\item[\textbf{2.61}] $S_k\le1+\sum_{n=2}^k\big(\tfrac1{n-1}-\tfrac1n\big)=2-\tfrac1k$.
\item[\textbf{2.62}] $c_n-c_{n+1}=\ln\tfrac{n+1}n-\tfrac1{n+1}>0$; e $c_n\ge\ln(n+1)-\ln n>0$ pelo teste da integral.
\item[\textbf{2.63}] $q!\,e-q!\sum_{n\le q}\tfrac1{n!}$ seria inteiro positivo menor que $\tfrac1{q+1}+\tfrac1{(q+1)^2}+\dots=\tfrac1q\le1$.
\item[\textbf{2.64}] Agrupe em blocos de $2^k$ termos; para $a_n=n^{-p}$, $\sum(2^{1-p})^k$ é geométrica.
\item[\textbf{2.65}] $a_n<1$ a partir de algum $n$, logo $a_n^2<a_n$; exemplo $a_n=\tfrac1{n^2}$.
\end{itemize}\end{multicols}}
